\section{Changing the speed without losing computation}\label{bcs:clocks-section}
A slower flow executes the same instructions only if its new clock reaches every finite logical time. The unforced viscous continuation in~\cite[Section~6.B]{CardonaMirandaPeraltaSalasPresas2021} uses a speed $Me^{-\nu t}$, with $M>0$, and has only the finite total logical time $M/\nu$. Here a prescribed residual force allows the nonintegrable speed $(1+t)^{-1}$. The logarithmic estimate below is the only clock bound needed for the first compact construction. Other onto clocks and spatial startup profiles follow that complete application in Section~\ref{bcs:clock-options}. The force is always defined from the prescribed field; it is not inferred from an unknown solution.

We use the explicit smooth step
\begin{equation}\label{bcs:smooth-step}
 \rho(v)=\begin{cases}e^{-1/v},&v>0,\\0,&v\le0,\end{cases}
 \qquad \sigma(v)=\frac{\rho(v)}{\rho(v)+\rho(1-v)}.
\end{equation}
It is zero for $v\le0$, one for $v\ge1$, and satisfies
$\sigma(v)+\sigma(1-v)=1$. Each positive-side derivative of $\rho$ is
$e^{-1/v}$ times a polynomial in $1/v$ and extends flatly by zero.
The bounds $y^m e^{-y}\le(m+k)!y^{-k}$ give effective errors near the
joining point, while the denominator of $\sigma$ is at least $e^{-2}$.
Thus all these functions and derivatives can be evaluated without an
exact equality test at a seam. For phase intervals we use
$\zeta(s)=\sigma(2s-1/2)$, which is constant near both endpoints.

\begin{lemma}[Onto changes of clock]\label{bcs:clock}
Fix $\nu>0$. Let $W(s,x)$ be a smooth divergence-free field for $s\ge0$, on a flat torus or on $\mathbb R^3$. Suppose all its mixed derivatives are bounded, $W(0)=0$, and in the whole-space case its spatial support lies in a fixed compact set. Let $g:[0,\infty)\to[0,\infty)$ be smooth, nondecreasing, onto, with $g(0)=0$. Then
\begin{align}
 u(t,x)&=g'(t)W(g(t),x),\\
 f(t,x)&=g''(t)W(g(t),x)-\nu g'(t)\Delta W(g(t),x)\notag\\
 &\quad +(g'(t))^2[\partial_sW+(W\cdot\nabla)W](g(t),x)
 \label{bcs:clock-force}
\end{align}
gives a prescribed solution with pressure zero and zero initial velocity. Its material flow is $X_u(t,a)=X_W(g(t),a)$, and hence preserves every all-time reachability event. If $(W\cdot\nabla)W=0$ identically, $f$ is divergence free. Spatial mean zero, when imposed on $W$, is inherited by $u$ and $f$.

For $g(t)=\log(1+t)$, every $j\ge0$ and spatial multi-index $\alpha$ satisfy
\begin{equation}\label{bcs:log-bounds}
 \|\partial_t^jD_x^\alpha u(t)\|_\infty+
 \|\partial_t^jD_x^\alpha f(t)\|_\infty
 \le C_{j,\alpha}(1+t)^{-1-j}.
\end{equation}
All these derivatives belong to $L^2([0,\infty);L^\infty)$; on a torus or with fixed compact support they also belong to space-time $L^2$.
\end{lemma}
\begin{proof}
The chain rule gives \eqref{bcs:clock-force} as $u_t+(u\cdot\nabla)u-\nu\Delta u$. The divergence statements follow by differentiation; the convection term is absent under the additional shear hypothesis. On a torus its spatial integral is zero because $(W\cdot\nabla)W=\operatorname{div}(W\otimes W)$. The same chain rule applied to $X_W(g(t),a)$ proves the trajectory identity. Onto means that every finite logical time has a finite preimage. This proves both directions of the event equivalence, including a visit during initialization. Existence is supplied by the explicit field, and uniqueness in the appropriate comparison class is Lemma~\ref{lem:b-comparison}.

For the logarithmic choice put $a=(1+t)^{-1}$ and $s=\log(1+t)$. For any smooth bounded-derivative function $A$ and positive integer $r$,
\[
 \partial_t[a^rA(s,x)]=a^{r+1}(\partial_s-r)A(s,x).
\]
The velocity is $aW$, while the force is
\[
 -\nu a\Delta W(s,x)+a^2[\partial_sW-W+(W\cdot\nabla)W](s,x).
\]
Repeated use of the displayed identity, commuting spatial derivatives, proves \eqref{bcs:log-bounds}. Integrating $(1+t)^{-2-2j}$ gives the stated norms. Fixed support also gives uniformly bounded kinetic energy. None of these calculations requires the executed instruction sequence.
\end{proof}
